% =====================================================================
% Appendices
% Each appendix begins with \chapter{...} (the \appendix command in
% main.tex switches \chapter numbering to A, B, C, ...).
% =====================================================================

\chapter{Declaration of Use of Generative AI}
\label{app:ai-declaration}

% ---------------------------------------------------------------------
% This appendix is REQUIRED. DTU permits the use of generative AI but
% mandates that students declare distinctly where and how it was used.
% The human authors remain the sole authors; AI is not a co-author.
% Sources: DTU library guidance and DTU AI rules (see CLAUDE.md §5).
% ---------------------------------------------------------------------

\section{Statement of Authorship and Responsibility}

We, the authors of this thesis, take full and sole responsibility for
the content of this report. All decisions regarding scope, claims,
results, and interpretations are ours. Generative artificial
intelligence (GAI) tools listed below were used as assistive
instruments only; they are not co-authors of this work.

\section{Tools Used}

\begin{table}[h]
\centering
\small
\begin{tabular}{@{}p{3.5cm}p{2.5cm}p{8cm}@{}}
\toprule
Tool & Version / Date & Purpose \\
\midrule
% Example row -- replace with actual usage before submission.
% Claude Code (Anthropic) & Opus 4.7, 2026 & Scaffolding LaTeX chapters; copy-editing; bibliography formatting. \\
\\
\bottomrule
\end{tabular}
\caption{Generative AI tools used during the preparation of this thesis.}
\label{tab:gai-tools}
\end{table}

\section{Where and How GAI Was Used}

% Fill one subsection per category of use. Be specific about which
% chapters / sections were affected. Vague disclosures ("used AI for
% writing") are not sufficient.

\subsection{Writing Assistance}
% E.g. grammar/style polishing, rephrasing for clarity. Identify the
% chapters or sections where this occurred. State that all factual
% claims and citations were verified by the author.

\subsection{Code Assistance}
% E.g. LaTeX scaffolding, build configuration, helper scripts in the
% reading-the-reader monorepo. State which files / commits.

\subsection{Literature Search and Synthesis}
% Only if applicable. If GAI was used to suggest search terms, state
% that no AI-generated citations were inserted without independent
% verification of the source.

\subsection{Figures, Tables, or Data}
% Only if applicable. If a figure or table was generated or
% restructured with AI, cite the tool in the figure/table caption
% as well as here.

\section{What GAI Was Not Used For}

% A non-exhaustive list of things the author did themselves, to make
% the boundary explicit. Example items below; revise to match reality.

\begin{itemize}
    \item Formulation of the problem statement and research questions.
    \item Selection and reading of all cited sources.
    \item Design decisions for the system architecture.
    \item Evaluation methodology and interpretation of results.
    \item Conclusions and discussion of limitations.
\end{itemize}

\section{Prompt and Output Retention}

Representative prompts and outputs are retained by the authors and
can be made available to the supervisor or examiners upon request,
in accordance with DTU guidance on documenting GAI use.

\section{Citations of GAI Output}

Where text, code, or figures derived from GAI output appear in the
body of the thesis, they are cited inline using the format below.

\begin{quote}
\small
Anthropic. \textit{Claude} (Opus 4.7) [Large Language Model].
\url{https://claude.ai}. Accessed YYYY-MM-DD.
\end{quote}

% =====================================================================
% Additional appendices below (extended results, full requirements
% list, raw evaluation data, repository structure, ethics approval,
% etc.). Add new chapters as needed.
% =====================================================================

% \chapter{Full Requirements Specification}
% \label{app:requirements}

% \chapter{Repository Structure}
% \label{app:repo}

% \chapter{Raw Evaluation Data}
% \label{app:raw-data}
